% GENERATED FILE. Edit canonical lessons, then run npm run generate:books.
% canonical-source-hash: fnv1a64:49d72151ad7fd04e
% canonical-lessons: HI-C375-asa, HI-C375-asanka, HI-C375-sanket, HI-C375-isara, HI-C375-rijhana

\chapter{Hopes and Fears}
\label{ch:hi-hopes-and-fears}
\hlchaptermodality{375}

\begin{chapteropening}
\textbf{By the end of this chapter:} \emph{I can say what I hope for and what I fear might happen.}

Complete the last lesson of chapter 375: I can say what I hope for and what I fear might happen.
\end{chapteropening}

\section[\texorpdfstring{\hi{आशा} (āśā)}{आशा (āśā)}]{\hi{आशा} (āśā) — a hope, an expectation}

\label{lesson:HI-C375-asa}



\begin{quote}
Before the new one: say the Hindi for usually, then the Hindi for to make out.
\end{quote}

\subsection*{You'll want to know: \hi{आशा}}
\textbf{\hi{आशा}} — \emph{āśā} — \textquotedblleft{}a hope, an expectation\textquotedblright{}.

\begin{quote}
\hi{मुझे} \hi{आशा} \hi{है} \hi{कि} \hi{अगर} \hi{मौसम} \hi{अच्छा} \hi{रहेगा} \hi{तो} \hi{फ़सल} \hi{अच्छी} \hi{होगी।}
\end{quote}

\begin{quote}
\emph{mujhe āśā hai ki agar mausam acchā rahegā to fasal acchī hogī.}
\end{quote}

\begin{quote}
I hope that if the weather stays good, the harvest will be good.
\end{quote}

Feminine.

\begin{cousinweb}
\textbf{\hi{आशा}} (\emph{āśā}) is a tatsama, borrowed from Sanskrit \emph{āśā}, \textquotedblleft{}a wish, a hope\textquotedblright{} (Monier-Williams).
\end{cousinweb}

\begin{grammarlens}[title={What you've built}]
One more word for saying what happens if.
\end{grammarlens}

\subsection*{Your turn}
\begin{itemize}
\raggedright
  \item \emph{Say it:} \emph{āśā}
  \item \emph{Say it:} \emph{āśā}, once more
  \item \emph{Say it:} \emph{mujhe āśā hai ki agar mausam acchā rahegā to fasal acchī hogī.} — I hope that if the weather stays good, the harvest will be good.
  \item \emph{Recall:} say the Hindi for usually, then the Hindi for to make out, then say \emph{āśā} again
\end{itemize}

\begin{tcolorbox}[breakable,colback=teal!4,colframe=teal!35!black,title={Before you move on}]
What does \hi{आशा} mean? (\textquotedblleft{}A hope, an expectation\textquotedblright{}.) Say it once more.
\end{tcolorbox}
\section[\texorpdfstring{\hi{आशंका} (āśaṅkā)}{आशंका (āśaṅkā)}]{\hi{आशंका} (āśaṅkā) — an apprehension, a fear (that something will happen)}

\label{lesson:HI-C375-asanka}



\begin{quote}
Before the new one: say the Hindi for to make out, then the Hindi for an expectation.
\end{quote}

\subsection*{You'll want to know: \hi{आशंका}}
\textbf{\hi{आशंका}} — \emph{āśaṅkā} — \textquotedblleft{}an apprehension, a fear (that something will happen)\textquotedblright{}.

\begin{quote}
\hi{आशंका} \hi{है} \hi{कि} \hi{अगर} \hi{बाढ़} \hi{आई} \hi{तो} \hi{पुल} \hi{टूट} \hi{जाएगा।}
\end{quote}

\begin{quote}
\emph{āśaṅkā hai ki agar bāṛh āī to pul ṭūṭ jāegā.}
\end{quote}

\begin{quote}
There are fears that if the flood comes, the bridge will collapse.
\end{quote}

Feminine.

\begin{cousinweb}
\textbf{\hi{आशंका}} (\emph{āśaṅkā}) is a tatsama, borrowed from Sanskrit \emph{āśaṅkā}, \textquotedblleft{}a doubt, a fear\textquotedblright{} (Monier-Williams).
\end{cousinweb}

\begin{grammarlens}[title={What you've built}]
One more word for saying what happens if.
\end{grammarlens}

\subsection*{Your turn}
\begin{itemize}
\raggedright
  \item \emph{Say it:} \emph{āśaṅkā}
  \item \emph{Say it:} \emph{āśaṅkā}, once more
  \item \emph{Say it:} \emph{āśaṅkā hai ki agar bāṛh āī to pul ṭūṭ jāegā.} — There are fears that if the flood comes, the bridge will collapse.
  \item \emph{Recall:} say the Hindi for to make out, then the Hindi for an expectation, then say \emph{āśaṅkā} again
\end{itemize}

\begin{tcolorbox}[breakable,colback=teal!4,colframe=teal!35!black,title={Before you move on}]
What does \hi{आशंका} mean? (\textquotedblleft{}An apprehension, a fear (that something will happen)\textquotedblright{}.) Say it once more.
\end{tcolorbox}
\section[\texorpdfstring{\hi{संकेत} (saṅket)}{संकेत (saṅket)}]{\hi{संकेत} (saṅket) — a sign, an indication}

\label{lesson:HI-C375-sanket}



\begin{quote}
Before the new one: say the Hindi for an expectation, then the Hindi for an apprehension.
\end{quote}

\subsection*{You'll want to know: \hi{संकेत}}
\textbf{\hi{संकेत}} — \emph{saṅket} — \textquotedblleft{}a sign, an indication\textquotedblright{}.

\begin{quote}
\hi{अगर} \hi{बच्चा} \hi{खाना} \hi{नहीं} \hi{खाता} \hi{तो} \hi{यह} \hi{बीमारी} \hi{का} \hi{संकेत} \hi{हो} \hi{सकता} \hi{है।}
\end{quote}

\begin{quote}
\emph{agar baccā khānā nahīṃ khātā to yah bīmārī kā saṅket ho saktā hai.}
\end{quote}

\begin{quote}
If a child won't eat, it can be a sign of illness.
\end{quote}

Masculine.

\begin{cousinweb}
\textbf{\hi{संकेत}} (\emph{saṅket}) is a tatsama, borrowed from Sanskrit \emph{saṃketa}, \textquotedblleft{}an agreement; a sign\textquotedblright{} (Monier-Williams).
\end{cousinweb}

\begin{grammarlens}[title={What you've built}]
One more word for saying what happens if.
\end{grammarlens}

\subsection*{Your turn}
\begin{itemize}
\raggedright
  \item \emph{Say it:} \emph{saṅket}
  \item \emph{Say it:} \emph{saṅket}, once more
  \item \emph{Say it:} \emph{agar baccā khānā nahīṃ khātā to yah bīmārī kā saṅket ho saktā hai.} — If a child won't eat, it can be a sign of illness.
  \item \emph{Recall:} say the Hindi for an expectation, then the Hindi for an apprehension, then say \emph{saṅket} again
\end{itemize}

\begin{tcolorbox}[breakable,colback=teal!4,colframe=teal!35!black,title={Before you move on}]
What does \hi{संकेत} mean? (\textquotedblleft{}A sign, an indication\textquotedblright{}.) Say it once more.
\end{tcolorbox}
\section[\texorpdfstring{\hi{इशारा} (iśārā)}{इशारा (iśārā)}]{\hi{इशारा} (iśārā) — a hint, a signal}

\label{lesson:HI-C375-isara}



\begin{quote}
Before the new one: say the Hindi for an apprehension, then the Hindi for an indication.
\end{quote}

\subsection*{You'll want to know: \hi{इशारा}}
\textbf{\hi{इशारा}} — \emph{iśārā} — \textquotedblleft{}a hint, a signal\textquotedblright{}.

\begin{quote}
\hi{अगर} \hi{मैं} \hi{इशारा} \hi{करूँगा} \hi{तो} \hi{अरुण} \hi{दरवाज़ा} \hi{खोलेगा।}
\end{quote}

\begin{quote}
\emph{agar maiṁ iśārā karūṅgā to Aruṇ darvāzā kholegā.}
\end{quote}

\begin{quote}
If I give a signal, Arun will open the door.
\end{quote}

Masculine.

\begin{cousinweb}
\textbf{\hi{इशारा}} (\emph{iśārā}) is usually traced to Arabic \emph{ishāra}, \textquotedblleft{}a sign, a gesture\textquotedblright{}, which reached Hindi through Persian.
\end{cousinweb}

\begin{grammarlens}[title={What you've built}]
One more word for saying what happens if.
\end{grammarlens}

\subsection*{Your turn}
\begin{itemize}
\raggedright
  \item \emph{Say it:} \emph{iśārā}
  \item \emph{Say it:} \emph{iśārā}, once more
  \item \emph{Say it:} \emph{agar maiṁ iśārā karūṅgā to Aruṇ darvāzā kholegā.} — If I give a signal, Arun will open the door.
  \item \emph{Recall:} say the Hindi for an apprehension, then the Hindi for an indication, then say \emph{iśārā} again
\end{itemize}

\begin{tcolorbox}[breakable,colback=teal!4,colframe=teal!35!black,title={Before you move on}]
What does \hi{इशारा} mean? (\textquotedblleft{}A hint, a signal\textquotedblright{}.) Say it once more.
\end{tcolorbox}
\section[\texorpdfstring{\hi{रिझाना} (rijhānā)}{रिझाना (rijhānā)}]{\hi{रिझाना} (rijhānā) — to charm, to win over}

\label{lesson:HI-C375-rijhana}



\begin{quote}
Before the new one: say the Hindi for an indication, then the Hindi for a hint.
\end{quote}

\subsection*{You'll want to know: \hi{रिझाना}}
\textbf{\hi{रिझाना}} — \emph{rijhānā} — \textquotedblleft{}to charm, to win over\textquotedblright{}.

\begin{quote}
\hi{अगर} \hi{तुम} \hi{अच्छा} \hi{गाना} \hi{गाओगे} \hi{तो} \hi{लोगों} \hi{को} \hi{रिझा} \hi{लोगे।}
\end{quote}

\begin{quote}
\emph{agar tum acchā gānā gāoge to logoṁ ko rijhā loge.}
\end{quote}

\begin{quote}
If you sing a good song, you'll win people over.
\end{quote}

Past, with \hi{ने} (he or she …): \hi{उसने} \hi{लोगों} \hi{को} \hi{रिझाया}, \emph{usne logoṁ ko rijhāyā}.

\begin{cousinweb}
\textbf{\hi{रिझाना}} (\emph{rijhānā}) is a tadbhava, inherited from Sanskrit \emph{ṛdhati}, \textquotedblleft{}thrives\textquotedblright{} (CDIAL 2457). CDIAL glosses the Hindi \emph{rijhānā} \textquotedblleft{}to enchant\textquotedblright{}.
\end{cousinweb}

\begin{grammarlens}[title={What you've built}]
One more word for saying what happens if.
\end{grammarlens}

\subsection*{Your turn}
\begin{itemize}
\raggedright
  \item \emph{Say it:} \emph{rijhānā}
  \item \emph{Say it:} \emph{rijhānā}, once more
  \item \emph{Say it:} \emph{agar tum acchā gānā gāoge to logoṁ ko rijhā loge.} — If you sing a good song, you'll win people over.
  \item \emph{Recall:} say the Hindi for an indication, then the Hindi for a hint, then say \emph{rijhānā} again
\end{itemize}

\begin{tcolorbox}[breakable,colback=teal!4,colframe=teal!35!black,title={Before you move on}]
What does \hi{रिझाना} mean? (\textquotedblleft{}To charm, to win over\textquotedblright{}.) Say it once more.
\end{tcolorbox}
